\section{楽器デバイスの同期プログラムのプログラムソース}
\label{sec:apndx_douki}
本節では，楽器デバイスの同期方式を担うArduinoプログラムの全文を示す．掲載にあたり，次の2点の省略を行っている．第一に，埋め込みの楽譜データ（Note[]配列）は，各小節ごとに同一形式のデータが多数繰り返されるのみであるため，冒頭2小節分のみを掲載し，以降は省略している．第二に，スレーブ機であるピアノ・ストリングス・フルートの3つのスケッチは，自機のIPアドレスと楽器ID（MY\_INST\_ID）の2定数および楽譜データを除いて完全に同一であるため，ピアノ（piano\_0706.ino）のみ全文を掲載し，ストリングス・フルートについては相違箇所のみを示す．

\subsection{drum\_0706.ino（マスター機）}
\begin{lstlisting}[caption={drum\_0706.ino（楽譜データは冒頭2小節のみ掲載）}, label={lst:douki_drum_full}]
#include <WiFiS3.h>
#include <WiFiUdp.h>
#include <Arduino.h>

const char SSID[] = "hackathon001-WPA2";
const char PASS[] = "hackathon001";

IPAddress localIP(192, 168, 11, 11); // ドラム（マスター）のIPアドレス
const uint16_t LOCAL_PORT = 5005; // 自身のポート番号

WiFiUDP Udp;
// UDP通信で受信したパケットのデータを一時的に格納する
uint8_t packetBuffer[32];

const int CDS_PIN = A0; // CdSセルを接続するアナログ入力ピン番号を定義
const int CDS_ON_THRESHOLD  = 800; // 起動直後だけ使う仮しきい値
const int CDS_OFF_THRESHOLD = 780; // 起動直後だけ使う仮しきい値
const int CDS_MIN_DELTA = 25;      // 暗い状態との差がこれ未満ならレーザー扱いしない
const uint8_t CDS_PEAK_HISTORY = 4; // 直近何回分のピークからしきい値を決めるか
const uint8_t CDS_ON_RATIO = 65;    // baselineからピークまでの65%をONしきい値にする
const uint8_t CDS_OFF_RATIO = 35;   // baselineからピークまでの35%をOFFしきい値にする
const int CDS_THRESHOLD_STEP_LIMIT = 25; // 1回の検出でしきい値が動きすぎない上限
const unsigned long CDS_STARTUP_CALIBRATION_MS = 1000; // 起動直後の誤検出防止時間
const unsigned long DEBOUNCE_MS = 20;// 直前の検知からこの時間未満の再検知を無視するデバウンス時間を定義
const int BPM_LIST[] = { 60, 90, 120, 150, 180 }; // 5段階の目標BPM設定値（60〜180）を定義
const float ROTATION_PERIOD_LIST[] = { 330, 240, 190, 150, 100 }; // 各BPMに対応する観覧車の実測回転周期を定義
const int MIN_STABLE_COUNT = 3; // BPM確定（安定）とみなす連続一致回数（3回）を定義

struct NoteData {
  uint16_t timing; // 発音タイミング（楽曲開始からの累積Tick数）
  uint16_t duration; // 発音の長さ（Tick単位）
  uint8_t  pitch; // 音高を示すMIDIノート番号（0〜127）
  uint8_t  startVelocity; // 発音開始時の強さ（0〜127）
  uint8_t  endVelocity; // 発音終了時の強さ（0〜127）
  uint8_t  type; // 発音の種類を指定するフラグ
};

// 楽譜データ
const NoteData Note[] = {

  // 【1小節目】
  {0, 1, 60, 80, 96, 0}, // [0] C4
  {12, 1, 64, 96, 80, 0}, // [1] E4
  {24, 1, 63, 80, 127, 0}, // [2] D#4
  {24, 1, 60, 80, 96, 0}, // [3] C4
  {36, 1, 64, 96, 80, 0}, // [4] E4
  {48, 1, 60, 80, 96, 0}, // [5] C4
  {60, 1, 64, 96, 80, 0}, // [6] E4
  {72, 1, 63, 80, 127, 0}, // [7] D#4
  {72, 1, 60, 80, 96, 0}, // [8] C4
  {84, 1, 64, 96, 80, 0}, // [9] E4

  // 【2小節目】
  {96, 1, 60, 80, 96, 0}, // [10] C4
  {108, 1, 64, 96, 80, 0}, // [11] E4
  {120, 1, 63, 80, 127, 0}, // [12] D#4
  {120, 1, 60, 80, 96, 0}, // [13] C4
  {132, 1, 64, 96, 80, 0}, // [14] E4
  {144, 1, 60, 80, 96, 0}, // [15] C4
  {156, 1, 64, 96, 80, 0}, // [16] E4
  {168, 1, 63, 80, 127, 0}, // [17] D#4
  {168, 1, 60, 80, 96, 0}, // [18] C4
  {180, 1, 64, 96, 80, 0}, // [19] E4
  {186, 1, 64, 48, 80, 0}, // [20] E4


  // …（以降，末尾まで同一形式の楽譜データが続くため省略）…

};

// NoteDataに含まれる音符の数
const uint16_t note_count = sizeof(Note) / sizeof(Note[0]);

// ==========================================
// Tick・演奏管理変数
// ==========================================
uint32_t globalTick     = 0;  // 全体基準Tick
uint32_t curTick        = 0;  // ドラム楽譜用Tick
uint32_t lastTick       = 0;  // updateTick用タイムスタンプ[µs]
uint32_t tickIval       = 0;  // 1Tickあたりのµs
uint16_t scoreIdx       = 0; // 次に発音する予定のデータのインデックス

uint8_t curBpm = 0; //現在のBPM
uint8_t curVol = 50; // 現在の音量

// ==========================================
// 楽譜シーケンス定義
//   startSequence[曲ID][イベント][0=演奏開始elapsedPlayTick, 1=楽器ID]
//   楽器ID: 0=ドラム, 1=ピアノ, 2=ストリングス, 3=フルート   終端={-1,-1}
//
//   ★ O通信送信タイミング = 演奏開始elapsedPlayTick - O_ADVANCE_TICKS
//   ★ state2移行後の最初のイベントがelapsedPlayTick=0（ドラム以外）の場合、
//     そのイベントのO通信は elapsedPlayTick=0 では間に合わないため、
//     実際の演奏開始は O_ADVANCE_TICKS 分だけ後ろにずれる
// ==========================================
const int32_t O_ADVANCE_TICKS = 96;  // O通信を演奏開始の何Tick前に送るか
const int     MAX_SONGS       = 5; // 最大の曲数
const int     MAX_EVENTS      = 10; // 最大の発音開始指示の上限
const int     drumInstId      = 0; // ドラムのID

// ★実際の曲構成に合わせて書き換える★
// elapsedPlayTick=0 からドラム以外の楽器が始まる場合、
// その楽器のtriggerTickを 0 ではなく O_ADVANCE_TICKS(=96) 以上に設定するか、
// もしくはドラム自身のtriggerTickを0にして、他楽器を96以降にする
int startSequence[MAX_SONGS][MAX_EVENTS][2] = {
  { // 曲ID: 0 かえるのうた
    {0,     drumInstId}, // elapsed=0でドラム自身が演奏開始
    {192,   1},          // elapsed=192でピアノ開始（O通信はelapsed=96）
    {384,   2},          // elapsed=384でストリングス（O通信はelapsed=288）
    {576,   3},          // elapsed=576でフルート（O通信はelapsed=480）
    {-1, -1}
  },
  { // 曲ID: 1 酒場でブギウギ
    //{0,     drumInstId}, // elapsed=0でドラム自身が演奏開始
    {0,   1},          // elapsed=192でピアノ開始（O通信はelapsed=96）
    //{0,   2},          // elapsed=384でストリングス（O通信はelapsed=288）
    //{0,   3},          // elapsed=576でフルート（O通信はelapsed=480）
    {-1, -1}
  },
  { // 曲ID: 2 王宮のメヌエット
    //{0,     drumInstId}, // elapsed=0でドラム自身が演奏開始
    //{0,   1},          // elapsed=192でピアノ開始（O通信はelapsed=96）
    {0,   2},          // elapsed=384でストリングス（O通信はelapsed=288）
    //{0,   3},          // elapsed=576でフルート（O通信はelapsed=480）
    {-1, -1}
  },
  { // 曲ID: 3 トルコ行進曲
    //{0,     drumInstId}, // elapsed=0でドラム自身が演奏開始
    {0,   1},          // elapsed=192でピアノ開始（O通信はelapsed=96）
    //{0,   2},          // elapsed=384でストリングス（O通信はelapsed=288）
    //{0,   3},          // elapsed=576でフルート（O通信はelapsed=480）
    {-1, -1}
  },
  { // 曲ID: 4 序曲
    {0,     drumInstId}, // elapsed=0でドラム自身が演奏開始
    {0,   1},          // elapsed=192でピアノ開始（O通信はelapsed=96）
    {0,   2},          // elapsed=384でストリングス（O通信はelapsed=288）
    {0,   3},          // elapsed=576でフルート（O通信はelapsed=480）
    {-1, -1}
  }
};

int     currentSongId       = 0; // 現在の曲のID
int32_t elapsedPlayTick     = 0;  // state3移行後（演奏開始後）の経過Tick

bool isDrumPlaying = false; //ドラムが演奏中でないかどうか

const uint8_t MAX_INSTRUMENTS = 4; // 最大楽器数
uint8_t oSentCount[MAX_EVENTS] = {0}; // 各イベントのO通信の送信回数を保存
bool isInstConnected[MAX_INSTRUMENTS]     = {false}; // 同期に参加してるかどうかv
bool isInstJoin[MAX_INSTRUMENTS]          = {false}; // 演奏同期に参加する予告が来たか
int32_t lastReceivedDiff[MAX_INSTRUMENTS] = {0}; // 他の楽器からT通信で送られるdiffを格納

uint8_t receivedBpmLevels[MAX_INSTRUMENTS] = {0}; // 参加している楽器のBPM段階を集約する配列（ドラム自身も含む）
uint8_t currentBpmLevel = 1; // 現在確定しているBPM段階（前回値）

// SyncManager（レーザー検知・BPM安定判定）
class SyncManager {
public:
  SyncManager()
    : lastDetectTime(0), correctedBPM(0),
      stableCount(0), prevCorrectedBPM(0), laserOn(false),
      baseline(0), baselineReady(false),
      onThreshold(CDS_ON_THRESHOLD), offThreshold(CDS_OFF_THRESHOLD),
      pulsePeak(0), peakIndex(0), peakCount(0),
      calibrated(false), startupCalibrated(false), waitingForRelease(false),
      calibrationStartMs(0), startupMin(1023), releaseThreshold(0) {
    clearPeakHistory();
  }

  // 全ての内部状態を初期値に戻す
  void reset() {
    lastDetectTime = 0;
    correctedBPM = 0;
    stableCount = 0;
    prevCorrectedBPM = 0;
    laserOn = false;
    baseline = 0;
    baselineReady = false;
    onThreshold = CDS_ON_THRESHOLD;
    offThreshold = CDS_OFF_THRESHOLD;
    pulsePeak = 0;
    peakIndex = 0;
    peakCount = 0;
    calibrated = false;
    startupCalibrated = false;
    waitingForRelease = false;
    calibrationStartMs = 0;
    startupMin = 1023;
    releaseThreshold = 0;
    clearPeakHistory();
  }

  // CdSの読み取り値と現在時刻から，しきい値較正・エッジ検出・デバウンスを行い，レーザー検知時にtrueを返す．
  bool update(int cdsValue, unsigned long nowMs) {
    if (calibrationStartMs == 0) calibrationStartMs = nowMs;
    if (!detectLight(cdsValue, nowMs)) return false;
    if (nowMs - lastDetectTime < DEBOUNCE_MS) return false;
    if (lastDetectTime != 0) {
      unsigned long interval = nowMs - lastDetectTime;
      int newBPM = intervalToBpm((float)interval);
      if (newBPM == prevCorrectedBPM) stableCount++;
      else                             stableCount = 1;
      prevCorrectedBPM = newBPM;
      correctedBPM     = newBPM;
    }
    lastDetectTime = nowMs;
    return true;
  }

  // 補正後の値を取得する窓口関数．
  int getCorrectedBPM() const { return correctedBPM; }
  int getStableCount()  const { return stableCount; }

private:
  unsigned long lastDetectTime; // 直前にレーザー検知が確定した時刻を格納
  int correctedBPM; // 今回レーザーが通過したときの間隔（BPM）
  int stableCount; // 同じBPMが何回連続で出たか
  int prevCorrectedBPM; // 前回レーザーが通過したときの間隔(BPM)
  bool laserOn; // 現在レーザーがCdSセルに当たっているかを格納
  int baseline; // 暗状態（レーザー非通過時）の読み取り値の追従値を格納
  bool baselineReady; // baselineの初期値取得が完了したかを格納
  int onThreshold; // 現在有効なON判定しきい値を格納
  int offThreshold; // 現在有効なOFF判定しきい値を格納
  int pulsePeak; // 検知中のパルスにおける読み取り値の最大値を一時的に格納
  int peakHistory[CDS_PEAK_HISTORY]; // しきい値較正に用いる直近ピーク値を保持する配列
  uint8_t peakIndex; // peakHistory内の現在の書き込み位置を管理するインデックス
  uint8_t peakCount; // peakHistoryに記録済みのピーク数を格納
  bool calibrated; // しきい値の較正が一度でも行われたかを格納
  bool startupCalibrated; // 起動直後のキャリブレーション期間が終了したかを格納
  bool waitingForRelease; // 起動時に暗転待ち状態かを格納
  unsigned long calibrationStartMs; // 起動キャリブレーションを開始した時刻を格納
  int startupMin; // 起動キャリブレーション中に観測された最小値を格納
  int releaseThreshold; // waitingForRelease状態を解除する判定値を格納

  // レーザーのON/OFF判定，起動時キャリブレーション，ベースライン更新を行う．
  bool detectLight(int v, unsigned long nowMs) {
    if (!baselineReady) {
      baseline = v;
      startupMin = v;
      baselineReady = true;
      updateFallbackThresholds();
    }

    // 起動直後にレーザーが当たっていても、最初の1秒は検出せず暗い側の値を探す
    if (!startupCalibrated) {
      if (v < startupMin) startupMin = v;
      baseline = startupMin;
      updateFallbackThresholds();
      if (nowMs - calibrationStartMs < CDS_STARTUP_CALIBRATION_MS) return false;
      startupCalibrated = true;
      waitingForRelease = (v >= onThreshold);
      releaseThreshold = (baseline > 1023 - CDS_MIN_DELTA) ? baseline - CDS_MIN_DELTA : offThreshold;
      return false;
    }

    if (waitingForRelease) {
      if (v <= releaseThreshold) {
        baseline = v;
        updateFallbackThresholds();
        waitingForRelease = false;
      }
      return false;
    }

    if (!laserOn) {
      if (v < baseline) baseline = v;
      else baseline = (baseline * 31 + v) / 32;
      if (!calibrated) updateFallbackThresholds();
    }

    if (!laserOn && v >= onThreshold) {
      laserOn = true;
      pulsePeak = v;
      return true;
    }

    if (laserOn) {
      if (v > pulsePeak) pulsePeak = v;
      if (v <= offThreshold) {
        laserOn = false;
        registerPeak(pulsePeak);
      }
    }
    return false;
  }

  // 検知したピーク値を記録し，ON/OFF判定しきい値を再較正する．
  void registerPeak(int peak) {
    if (peak - baseline < CDS_MIN_DELTA) return;

    peakHistory[peakIndex] = peak;
    peakIndex = (peakIndex + 1) % CDS_PEAK_HISTORY;
    if (peakCount < CDS_PEAK_HISTORY) peakCount++;
    if (peakCount < 3) return;

    long sum = 0;
    for (uint8_t i = 0; i < peakCount; i++) sum += peakHistory[i];
    int peakAvg = sum / peakCount;
    int amplitude = peakAvg - baseline;
    if (amplitude < CDS_MIN_DELTA) return;

    int targetOn = constrain(baseline + (amplitude * CDS_ON_RATIO) / 100, 0, 1023);
    int targetOff = constrain(baseline + (amplitude * CDS_OFF_RATIO) / 100, 0, targetOn - 5);

    if (!calibrated) {
      onThreshold = targetOn;
      offThreshold = targetOff;
      calibrated = true;
    } else {
      onThreshold = limitStep(onThreshold, targetOn, CDS_THRESHOLD_STEP_LIMIT);
      offThreshold = limitStep(offThreshold, targetOff, CDS_THRESHOLD_STEP_LIMIT);
    }
    if (offThreshold >= onThreshold) offThreshold = max(0, onThreshold - 5);
  }

  // 較正前に使用する仮のしきい値を更新する．
  void updateFallbackThresholds() {
    onThreshold = constrain(baseline + CDS_MIN_DELTA, 0, 1023);
    offThreshold = constrain(baseline + CDS_MIN_DELTA / 2, 0, max(0, onThreshold - 5));
  }

  // しきい値が一度に変化する最大量を制限する．
  int limitStep(int current, int target, int limit) {
    if (target > current + limit) return current + limit;
    if (target < current - limit) return current - limit;
    return target;
  }

  // ピーク履歴を初期化する．
  void clearPeakHistory() {
    for (uint8_t i = 0; i < CDS_PEAK_HISTORY; i++) peakHistory[i] = 0;
  }

  // 実測の1周時間テーブルの中から最も近いものを探し、対応するBPM_LISTの値を返す
  int intervalToBpm(float intervalMs) {
    int best = BPM_LIST[0]; float bestD = fabs(intervalMs - ROTATION_PERIOD_LIST[0]);
    for (int i = 1; i < 5; i++) {
      float d = fabs(intervalMs - ROTATION_PERIOD_LIST[i]);
      if (d < bestD) { bestD = d; best = BPM_LIST[i]; }
    }
    return best;
  }
};

SyncManager syncManager;

// ==========================================
// state:
//   0 = 待機（BPM安定待ち）
//   1 = BPM安定、楽器と同期中（diff収束待ち）
//   2 = 同期完了・演奏開始前（O通信準備 / state3への96tick待機）
//   3 = 演奏中（elapsedPlayTick が進む、ドラムも演奏する場合は isDrumPlaying=true）
//   4 = システム停止
// ==========================================
int state = 4;

// 楽譜遷移（ジャンプ）定義
// jumpTable[曲ID][ジャンプ番号][0=トリガーtick, 1=ジャンプ先tick, 2=ジャンプ先scoreIdx]
const int MAX_JUMPS = 20;
const int32_t jumpTable[MAX_SONGS][MAX_JUMPS][3] = {
  { // 曲ID: 0 かえるのうた
    {768, 0, 0},
    {-1, -1, -1}
  },
  {
    {768, 0, 0},
    {-1, -1, -1}
  },
  {
    {768, 0, 0},
    {-1, -1, -1}
  },
  {
    {768, 0, 0},
    {-1, -1, -1}
  },
  {
    {3624, 2088, 637},
    {-1, -1, -1}
  }
};

uint8_t currentJumpIdx = 0; // 現在のジャンプインデックス

// 曲ごとの初期設定データ
struct SongStartData {
  uint32_t startTick;
  uint16_t startNoteIdx;
};

// Processingの設定に合わせた初期値の定義
const SongStartData songStartParams[MAX_SONGS] = {
  {0, 0},   // 曲ID: 0 (かえるのうた)
  {0, 0},   // 曲ID: 1 (酒場でブギウギ風)
  {0, 0},   // 曲ID: 2 (王宮のメヌエット)
  {0, 0},    // 曲ID: 3 (トルコ行進曲)
  {768, 92} // 曲ID: 4 (序曲)
};





// 関数宣言
void updateTick();
void setBpm(uint8_t bpm);
void sendVolumeValue(uint8_t vol);
void sendPacket();
void processOCommands();
void sendOPacket(uint8_t instId, uint8_t songId, uint32_t startGlobalTick);
bool checkSyncReady();
void sendNPacket(uint8_t bpm);
void sendRPacket(uint8_t instId, int32_t diff);

// IPアドレス（楽器IDからIPを引く）
IPAddress instIP(uint8_t instId) {
  // ID=1→192.168.11.12, ID=2→.13, ID=3→.14
  return IPAddress(192, 168, 11, 11 + instId);
}

// セットアップ
void setup() {
  Serial.begin(115200);
  WiFi.config(localIP);
  WiFi.begin(SSID, PASS);
  Serial.print("Connecting WiFi...");
  while (WiFi.status() != WL_CONNECTED) { delay(500); Serial.print("."); }
  Serial.println("\nWiFi Connected! IP: " + WiFi.localIP().toString());
  Udp.begin(LOCAL_PORT);
  Serial.println("state=0 BPM安定待ち");
}

// メインループ
void loop() {
  unsigned long now      = millis();
  int           cdsValue = analogRead(CDS_PIN);

  // ---- UDP受信 ----
  int pktSize = Udp.parsePacket();
  while (pktSize > 0) {
  if (pktSize >= 2) {
    Udp.read(packetBuffer, sizeof(packetBuffer));
    char header = (char)packetBuffer[0];

    switch (header) {
      //楽器からの同期パケット
      case 'T':
        if (pktSize >= 6) {
          uint32_t slaveTick =
            ((uint32_t)packetBuffer[1] << 24) |
            ((uint32_t)packetBuffer[2] << 16) |
            ((uint32_t)packetBuffer[3] << 8)  |
             (uint32_t)packetBuffer[4];
          uint8_t id = packetBuffer[5];

          if (id < MAX_INSTRUMENTS && id != drumInstId) {
            isInstConnected[id] = true;
            if (isInstJoin[id] != true) isInstJoin[id] = true;
            int32_t diff = (int32_t)(globalTick - slaveTick);
            lastReceivedDiff[id] = diff;

            sendRPacket(id, diff);

            Serial.print(F("[T] id=")); Serial.print(id);
            Serial.print(F(" diff=")); Serial.println(diff);
          }
        }
        break;

      case 'V': {
        uint8_t step = packetBuffer[1];
        uint8_t vol  = (step >= 1 && step <= 5) ? (step * 10) : 30;
        sendVolumeValue(vol);
        break;
      }
      //指揮機からのブロードキャスト終了指示．演奏を停止し、state4にして初期化する
      case 'E':
        isDrumPlaying       = false;
        curTick             = 0;
        scoreIdx            = 0;
        globalTick          = 0;
        elapsedPlayTick     = 0;
        state               = 4;
        memset(oSentCount,         0, sizeof(oSentCount));
        memset(isInstConnected,    0, sizeof(isInstConnected));
        memset(isInstJoin,         0, sizeof(isInstJoin));
        curBpm = 0;
        tickIval = 0;

        // BPM関連の初期化
        memset(receivedBpmLevels, 0, sizeof(receivedBpmLevels));
        currentBpmLevel = 0;

        Serial.println(F("[E] 演奏終了。state=4 に初期化"));
        break;

      // 指揮機からのブロードキャスト開始指示．state0に戻して初期化する
      case 'S': 
      {
        static unsigned long lastSReceivedTime = 0;
        unsigned long nowMs = millis();
        
        // 1秒以内の連続したS通信（UDP連送）は無視する
        if (nowMs - lastSReceivedTime < 1000) {
          Serial.println(F("[S] 連続受信のためスキップ"));
          break;
        }
        lastSReceivedTime = nowMs;

        // 楽譜IDのローテーション
        currentSongId = (currentSongId + 1) % MAX_SONGS;

        // 初期化処理
        isDrumPlaying       = false;
        curTick             = 0;
        scoreIdx            = 0;
        globalTick          = 0;
        elapsedPlayTick     = 0;
        currentJumpIdx      = 0; // ジャンプインデックスも初期化
        state               = 0;
        memset(oSentCount,         0, sizeof(oSentCount));
        memset(isInstConnected,    0, sizeof(isInstConnected));
        memset(isInstJoin,         0, sizeof(isInstJoin));
        curBpm = 0;
        tickIval = 0;

        memset(receivedBpmLevels, 0, sizeof(receivedBpmLevels));
        currentBpmLevel = 0;

        syncManager.reset();

        Serial.print(F("[S] 演奏開始。state=0 に初期化。曲変更。SongID:"));
        Serial.println(currentSongId);
        break;
      }

      // 'J': 楽器からの演奏参加予告．同期に参加する楽器を読み取る
      case 'J':
      if(pktSize >= 2){
        uint8_t instId = packetBuffer[1];
        if (instId < MAX_INSTRUMENTS) {
          isInstJoin[instId] = true;
          Serial.print(F("[J] 参加予告受信: instId=")); Serial.println(instId);
        }
      }
      break;

      // 子機からのBPM提案を受信
      case 'M':
        if (pktSize >= 3) {
          uint8_t bpmVal = packetBuffer[1];
          uint8_t instId = packetBuffer[2];
          if (instId < MAX_INSTRUMENTS) {
            receivedBpmLevels[instId] = bpmVal;
            Serial.print(F("[M] BPM提案受信 instId=")); Serial.print(instId);
            Serial.print(F(" BPM=")); Serial.println(bpmVal);
          }
        }
        
        break;
        

      default: break;
    }
  } else {
    Udp.read(); // 空読みして廃棄
  }
  pktSize = Udp.parsePacket(); // 次のパケットがあるか確認
  }

  if (state == 4) return;

  // ---- レーザー検知 ----
  bool laser = syncManager.update(cdsValue, now);

  // ---- ステートマシン ----
  switch (state) {

    // --- state0: BPM安定待ち ---
    case 0:
      if (laser) {
        int s = syncManager.getStableCount();
        uint8_t curLaserBpm = syncManager.getCorrectedBPM();
        Serial.print(F("[S0] BPM=")); Serial.print(curLaserBpm);
        Serial.print(F(" stable=")); Serial.println(s);

        // ドラム自身のレーザーBPMを配列に上書き
        receivedBpmLevels[drumInstId] = curLaserBpm;

        if (s >= MIN_STABLE_COUNT) {
          state = 1;
          lastTick = micros();

          // 初回確定時に自身へ適用 ＆ 即座にN通信で子機に送信
          currentBpmLevel = curLaserBpm;
          setBpm(currentBpmLevel);
          sendNPacket(currentBpmLevel);

          Serial.println(F("[S0→1] BPM安定"));
        }
      }
      break;

    // --- state1: 楽器と同期待ち ---
    case 1:
      updateTick();
      if (laser) {
        receivedBpmLevels[drumInstId] = syncManager.getCorrectedBPM(); // 配列上書きのみ
        if (checkSyncReady()) {
          // state2へ移行。elapsedPlayTick は state3移行後に開始するが、O通信の送信タイミング算出のためにカウントはここから始める
          state = 2;
          elapsedPlayTick = -O_ADVANCE_TICKS;
          memset(oSentCount, 0, sizeof(oSentCount));
          Serial.println(F("[S1→2] 同期完了"));
        } else {
          Serial.println(F("[S1] 同期待ち..."));
        }
      }
      break;

    // --- state2: 同期完了・演奏開始待ち ---
    // elapsedPlayTick を進めながら O通信タイミングを監視する。state3への移行は processOCommands() 内でドラム自身の開始タイミングが来たとき。
    case 2:
      updateTick();
      if (laser) receivedBpmLevels[drumInstId] = syncManager.getCorrectedBPM();
      processOCommands();
      break;

    // --- state3: 演奏中 ---
    case 3:
      updateTick();
      if (laser) receivedBpmLevels[drumInstId] = syncManager.getCorrectedBPM();
      processOCommands();
      break;
    
    // --- state4: システム停止中 ---
    // case 4: break;
  }
}

// 同期完了チェック
//   接続済み楽器が1台以上 かつ 全員のdiffが SYNC_READY_THRESHOLD以内
bool checkSyncReady() {
  bool expectedToJoin = false;
  
  for (uint8_t i = 1; i < MAX_INSTRUMENTS; i++) {
    if (isInstJoin[i]) {
      expectedToJoin = true; // 参加予定の楽器が少なくとも1台いる

      float SYNC_READY_THRESHOLD = 60 * 1000 / tickIval;// 60ms超えてたら同期ができていないとみなすから50
      
      // 参加予定だが未接続、または同期のズレが大きい場合は「まだ準備未完了」
      if (!isInstConnected[i] || abs(lastReceivedDiff[i]) > SYNC_READY_THRESHOLD) {
        return false;
      }
    }
  }
  
  // 参加予定の全楽器が条件をクリアしていれば true。
  // ※もし「他楽器が0台（ドラムソロ）でも開始させたい」場合は、return true;
  return expectedToJoin; 
}

// ==========================================
// O通信管理
//   各イベントについて「演奏開始elapsedTick - O_ADVANCE_TICKS」になったら送信。
//   ドラム自身のイベント（ID=0）は演奏開始処理のみ行う（O通信不要）。
//
//   ・elapsedPlayTick が負（=state2の最初96tick）の間も処理するため、
//     startSequenceの最初のイベントが elapsed=0 であっても対応できる。
//   ・受信側(notDrum)が演奏開始予定globalTick をもとにcurTickを計算するため、
//     「ドラムの globalTick + O_ADVANCE_TICKS」を startGlobalTick として送る。
// ==========================================
void processOCommands() {
  if (currentSongId < 0 || currentSongId >= MAX_SONGS) return;

  for (int i = 0; i < MAX_EVENTS; i++) {
    int32_t evTick = startSequence[currentSongId][i][0];
    int     instId = startSequence[currentSongId][i][1];
    
    if (evTick == -1) break;
    if (oSentCount[i] >= 3) continue; // 3回送信し終わっていたらスキップ

    int32_t ticksUntilStart = evTick - elapsedPlayTick;

    if (instId == drumInstId) {
      if (elapsedPlayTick >= evTick) {
        if (!isDrumPlaying) {
          isDrumPlaying = true;
          curTick  = songStartParams[currentSongId].startTick;
          scoreIdx = songStartParams[currentSongId].startNoteIdx;
          state    = 3;
          Serial.println(F("[DRUM] 演奏開始"));
        }
        oSentCount[i] = 3; // ドラム自身は通信不要なので完了扱い
      }
    } else {
      if (isInstJoin[instId] && isInstConnected[instId]) {
        uint32_t startGlobalTick = (uint32_t)((int32_t)globalTick + ticksUntilStart);

        // O_ADVANCE_TICKS(96)を切っていて、まだ0回目なら 1回目の送信
        if (ticksUntilStart <= O_ADVANCE_TICKS && oSentCount[i] == 0) {
          sendOPacket_Single(instId, currentSongId, startGlobalTick);
          oSentCount[i] = 1;
        } 
        // 2/3 (64Tick前) を切っていて、まだ1回目なら 2回目の送信
        else if (ticksUntilStart <= O_ADVANCE_TICKS * 2 / 3 && oSentCount[i] == 1) {
          sendOPacket_Single(instId, currentSongId, startGlobalTick);
          oSentCount[i] = 2;
        } 
        // 1/3 (32Tick前) を切っていて、まだ2回目なら 3回目の送信
        else if (ticksUntilStart <= O_ADVANCE_TICKS / 3 && oSentCount[i] == 2) {
          sendOPacket_Single(instId, currentSongId, startGlobalTick);
          oSentCount[i] = 3;
        }
      }
    }
  }
  if (state == 2 && !isDrumPlaying) {
    int32_t earliestEvent = INT32_MAX;
    for (int i = 0; i < MAX_EVENTS; i++) {
      int32_t evTick = startSequence[currentSongId][i][0];
      int     instId = startSequence[currentSongId][i][1];
      if (evTick == -1) break;
      
      // 修正点：ドラム自身（drumInstId）も開始時間の計算に含める
      if ((instId == drumInstId || isInstJoin[instId]) && evTick < earliestEvent) {
        earliestEvent = evTick;
      }
    }

    if (earliestEvent == INT32_MAX || earliestEvent > O_ADVANCE_TICKS * 2) {
      if (elapsedPlayTick < 0) {
        elapsedPlayTick = 0;
        Serial.println(F("[SKIP] 近接楽器なし: elapsedPlayTick=0 に即進め"));
      }
    } else {
      int32_t targetElapsed = earliestEvent - O_ADVANCE_TICKS;
      if (elapsedPlayTick < targetElapsed) {
        elapsedPlayTick = targetElapsed;
        Serial.print(F("[SKIP] elapsedPlayTick を ")); Serial.print(targetElapsed);
        Serial.println(F(" に調整"));
      }
    }
  }
}

// O通信を「1回だけ(delayなしで)」送る専用の関数
void sendOPacket_Single(uint8_t instId, uint8_t songId, uint32_t startGlobalTick) {
  Udp.beginPacket(instIP(instId), LOCAL_PORT);
  Udp.write('O');
  Udp.write(songId);
  Udp.write((uint8_t)((startGlobalTick >> 24) & 0xFF));
  Udp.write((uint8_t)((startGlobalTick >> 16) & 0xFF));
  Udp.write((uint8_t)((startGlobalTick >>  8) & 0xFF));
  Udp.write((uint8_t)( startGlobalTick        & 0xFF));
  Udp.endPacket();
}

// N通信パケット送信 (確定BPMを全参加楽器へブロードキャスト)
void sendNPacket(uint8_t bpm) {
    for (int i = 1; i < MAX_INSTRUMENTS; i++) {
      Udp.beginPacket(instIP(i), LOCAL_PORT);
      Udp.write('N');
      Udp.write(bpm);
      Udp.endPacket();
      // }
      delay(1);//2msならok
    }
  // }
}

void sendRPacket(uint8_t instId,int32_t diff){
  // R通信: diffのみ返信（O通信は別途 sendOPacket で送る）
  Udp.beginPacket(instIP(instId), LOCAL_PORT);
  Udp.write('R');
  Udp.write((uint8_t)((diff >> 24) & 0xFF));
  Udp.write((uint8_t)((diff >> 16) & 0xFF));
  Udp.write((uint8_t)((diff >>  8) & 0xFF));
  Udp.write((uint8_t)( diff        & 0xFF));
  Udp.endPacket();
}

// ==========================================
// Tick管理
// ==========================================
void updateTick() {
  if (tickIval == 0) return;
  uint32_t now = micros();
  while (now - lastTick >= tickIval) {
    uint32_t currentIval = tickIval; // ループ突入時の値を保存

    globalTick++;

    // --- 48TickごとにBPM段階の多数決とN通信送信 ---
    if (state >= 1 && globalTick % 48 == 0) {
      uint8_t newLevel = determineSystemBpmLevel();
      if (newLevel != currentBpmLevel && newLevel > 0) {
        currentBpmLevel = newLevel;
        setBpm(currentBpmLevel); 
      }
      sendNPacket(currentBpmLevel);
      // Serial.print(F("[SYNC] 48Tick周期 N通信送信: 確定BPM = ")); 
      // Serial.println(currentBpmLevel);
    }

    if (state >= 2) elapsedPlayTick++;
    if (state == 3 && isDrumPlaying) {
      if (jumpTable[currentSongId][currentJumpIdx][0] != -1 && curTick == jumpTable[currentSongId][currentJumpIdx][0]) {
        curTick  = jumpTable[currentSongId][currentJumpIdx][1];
        scoreIdx = jumpTable[currentSongId][currentJumpIdx][2];
        currentJumpIdx++;
        
        if (currentJumpIdx >= MAX_JUMPS || jumpTable[currentSongId][currentJumpIdx][0] == -1) {
          currentJumpIdx = 0;
        }
        Serial.println(F("[JUMP] 楽譜をジャンプしました！"));
      }

      sendPacket();
      curTick++;
    }
    
    lastTick += currentIval; // 書き換えられていても元の値を足す
  }
}

// ==========================================
// BPM更新 → tickIval再計算
// ==========================================
void setBpm(uint8_t bpm) {
  if (bpm == curBpm) return;
  curBpm   = bpm;
  tickIval = 60000000UL / (uint32_t)bpm / 24;
  Serial.print('B'); Serial.print(','); Serial.print(bpm); Serial.print('\n');
}

void sendVolumeValue(uint8_t vol) {
  if (vol == curVol) return;
  curVol = vol;
  Serial.print('V'); Serial.print(','); Serial.print(vol); Serial.print('\n');
}

// N通信で送信するためのBPM段階を決定する関数
uint8_t determineSystemBpmLevel() {
  uint8_t validLevels[MAX_INSTRUMENTS];
  uint8_t validCount = 0;

  // 有効なデータ（参加中かつデータがある楽器）だけを抽出
  for (int i = 0; i < MAX_INSTRUMENTS; i++) {
    if ((i == drumInstId || isInstJoin[i]) && receivedBpmLevels[i] > 0) {
      validLevels[validCount] = receivedBpmLevels[i];
      validCount++;
    }
  }

  if (validCount == 0) return currentBpmLevel; // データが一つもなければ前回維持
  if (validCount == 1) return validLevels[0];  // 1台だけならそれを採用

  // 2. 配列を小さい順にバブルソート
  for (int i = 0; i < validCount - 1; i++) {
    for (int j = 0; j < validCount - i - 1; j++) {
      if (validLevels[j] > validLevels[j + 1]) {
        uint8_t temp = validLevels[j];
        validLevels[j] = validLevels[j + 1];
        validLevels[j + 1] = temp;
      }
    }
  }

  // 3. 中央値の取得と、偶数時の「ユーザー提案ロジック（丸め）」
  if (validCount % 2 != 0) {
    // 奇数個なら、ど真ん中がそのまま中央値
    return validLevels[validCount / 2];
  } else {
    // 偶数個なら、中央を挟む2つの値を取得
    uint8_t mid1 = validLevels[validCount / 2 - 1];
    uint8_t mid2 = validLevels[validCount / 2];

    // もし2つの値が同じなら、それが中央値
    if (mid1 == mid2) return mid1;

    // 異なる場合、前回値（currentBpmLevel）と差が小さい方を採用
    int diff1 = abs((int)mid1 - (int)currentBpmLevel);
    int diff2 = abs((int)mid2 - (int)currentBpmLevel);

    if (diff1 <= diff2) {
      return mid1;
    } else {
      return mid2;
    }
  }
}

// ==========================================
// ノート送信（Processingへ）
// ==========================================
void sendPacket() {
  while (scoreIdx < note_count && Note[scoreIdx].timing <= curTick) {
    //uint8_t vs = (uint8_t)((uint16_t)Note[scoreIdx].startVelocity * curVol / 50);
    //uint8_t ve = (uint8_t)((uint16_t)Note[scoreIdx].endVelocity   * curVol / 50);
    Serial.print('N'); Serial.print(',');
    Serial.print(Note[scoreIdx].pitch);    Serial.print(',');
    Serial.print(Note[scoreIdx].duration); Serial.print(',');
    Serial.print(Note[scoreIdx].startVelocity); Serial.print(',');
    Serial.print(Note[scoreIdx].endVelocity); Serial.print(',');
    Serial.println(Note[scoreIdx].type);
    scoreIdx++;
  }
}
\end{lstlisting}

\subsection{piano\_0706.ino（スレーブ機）}
\begin{lstlisting}[caption={piano\_0706.ino（楽譜データは冒頭2小節のみ掲載）}, label={lst:douki_piano_full}]
#include <WiFiS3.h>
#include <WiFiUdp.h>

const char SSID[] = "hackathon001-WPA2";
const char PASS[] = "hackathon001";

// ★楽器ごとに変更（ピアノ=12, ストリングス=13, フルート=14）
IPAddress localIP(192, 168, 11, 12);
IPAddress drumIP (192, 168, 11, 11);
const uint16_t LOCAL_PORT = 5005;

// ★楽器ID（0=ドラム, 1=ピアノ, 2=ストリングス, 3=フルート）
const uint8_t MY_INST_ID = 1;

WiFiUDP Udp;
// UDP通信で受信したパケットのデータを一時的に格納
uint8_t packetBuffer[32];

const int CDS_PIN           = A0; //CdSセルを接続するアナログ入力ピン番号を定義
const int CDS_ON_THRESHOLD  = 820; // 起動直後だけ使う仮しきい値
const int CDS_OFF_THRESHOLD = 800;  // 起動直後だけ使う仮しきい値
const int CDS_MIN_DELTA = 25;       // 暗い状態との差がこれ未満ならレーザー扱いしない
const uint8_t CDS_PEAK_HISTORY = 4; // 直近何回分のピークからしきい値を決めるか
const uint8_t CDS_ON_RATIO = 65;    // baselineからピークまでの65%をONしきい値にする
const uint8_t CDS_OFF_RATIO = 35;   // baselineからピークまでの35%をOFFしきい値にする
const int CDS_THRESHOLD_STEP_LIMIT = 25; // 1回の検出でしきい値が動きすぎない上限
const unsigned long CDS_STARTUP_CALIBRATION_MS = 1000; // 起動直後の誤検出防止時間
const unsigned long DEBOUNCE_MS = 20;// 直前の検知からこの時間未満の再検知を無視するデバウンス時間を定義
const int BPM_LIST[] = { 60, 90, 120, 150, 180 };
// 段階1~5(BPM_LISTと対応)の実測の1周あたり時間[ms]。理論値(30000/BPM)では観覧車の
// 実際の回転速度と合わないことが判明したため、実測値でBPMを判定する。
const float ROTATION_PERIOD_LIST[] = {380, 280, 210, 160, 100};//pwmと大体対応させている
const int MIN_STABLE_COUNT  = 3; //BPM確定（安定）とみなす連続一致回数を定義

// 楽譜データ
struct NoteData {
  uint16_t timing; // 発音タイミング
  uint16_t duration; // 発音の長さ（Tick単位）
  uint8_t  pitch; // 音高を示すMIDIノート番号（0〜127）
  uint8_t  startVelocity; // 発音開始時の強さ（0〜127）
  uint8_t  endVelocity; // 発音開始時の強さ（0〜127）
  uint8_t  type; // 発音開始時の強さ（0〜127）
};

const NoteData Note[] = {

  // 【1小節目】
  {0, 24, 60, 96, 80, 0}, // [0] C4
  {24, 24, 62, 96, 80, 0}, // [1] D4
  {48, 24, 64, 96, 80, 0}, // [2] E4
  {72, 24, 65, 96, 80, 0}, // [3] F4

  // 【2小節目】
  {96, 24, 64, 96, 80, 0}, // [4] E4
  {120, 24, 62, 96, 80, 0}, // [5] D4
  {144, 36, 60, 96, 80, 0}, // [6] C4


  // …（以降，末尾まで同一形式の楽譜データが続くため省略）…

};

const uint16_t note_count = sizeof(Note) / sizeof(Note[0]); //音符の数(配列に含まれている合計/音符1つあたりの要素数)
uint16_t scoreIdx = 0; // 次に発音する予定のデータのインデックス

// ==========================================
// Tick管理変数
// ==========================================
uint32_t globalTick   = 0; //全体基準のTick
uint32_t curTick      = 0;  // 楽譜用Tick
uint32_t lastTick     = 0; // updateTick用タイムスタンプ

uint32_t baseTickIval = 0;   // BPMから計算した基本interval [µs/tick]
uint32_t tickIval     = 0;   // 実際に使うinterval（同期補正あり）
int32_t  ticksToCorrect = 0; // 補正残りTick数
int32_t  lastDiff       = 0; // 直近のdiff（BPM変化時の再計算用）
uint32_t lastTPacketMicroTime = 0; //直近のTパケット送信時の時刻Micro
uint32_t lastTPacketTick = 0; // 直近のTパケット送信時のglobalTick（drumNotStarted判定用）

bool drumNotStarted = true;// ドラム起動前同期フラグ

uint8_t curBpm = 0; // 現在のBPM
uint8_t curVol = 50; // 現在の音量

// 演奏管理
bool     waitingForStart  = false; // O通信を受信して演奏開始globalTickを待っている
uint32_t startGlobalTick  = 0;  // ドラムから指定された演奏開始globalTick
uint8_t  currentSongId    = 0; // 現在の曲のID

// SyncManager（レーザー検知・BPM安定判定）
class SyncManager {
public:
  SyncManager()
    : lastDetectTime(0), correctedBPM(0),
      stableCount(0), prevCorrectedBPM(0), laserOn(false),
      baseline(0), baselineReady(false),
      onThreshold(CDS_ON_THRESHOLD), offThreshold(CDS_OFF_THRESHOLD),
      pulsePeak(0), peakIndex(0), peakCount(0),
      calibrated(false), startupCalibrated(false), waitingForRelease(false),
      calibrationStartMs(0), startupMin(1023), releaseThreshold(0) {
    clearPeakHistory();
  }

  // 全ての内部状態を初期値に戻す
  void reset() {
    lastDetectTime = 0;
    correctedBPM = 0;
    stableCount = 0;
    prevCorrectedBPM = 0;
    laserOn = false;
    baseline = 0;
    baselineReady = false;
    onThreshold = CDS_ON_THRESHOLD;
    offThreshold = CDS_OFF_THRESHOLD;
    pulsePeak = 0;
    peakIndex = 0;
    peakCount = 0;
    calibrated = false;
    startupCalibrated = false;
    waitingForRelease = false;
    calibrationStartMs = 0;
    startupMin = 1023;
    releaseThreshold = 0;
    clearPeakHistory();
  }

  // CdSの読み取り値と現在時刻から，しきい値較正・エッジ検出・デバウンスを行い，レーザー検知時にtrueを返す．
  bool update(int cdsValue, unsigned long nowMs) {
    if (calibrationStartMs == 0) calibrationStartMs = nowMs;
    if (!detectLight(cdsValue, nowMs)) return false;
    if (nowMs - lastDetectTime < DEBOUNCE_MS) return false;
    if (lastDetectTime != 0) {
      unsigned long interval = nowMs - lastDetectTime;
      int newBPM = intervalToBpm((float)interval);
      if (newBPM == prevCorrectedBPM) stableCount++;
      else                             stableCount = 1;
      prevCorrectedBPM = newBPM;
      correctedBPM     = newBPM;
    }
    lastDetectTime = nowMs;
    return true;
  }

  // 補正後の値を取得する窓口関数．
  int getCorrectedBPM() const { return correctedBPM; }
  int getStableCount()  const { return stableCount; }

private:
  unsigned long lastDetectTime; // 直前にレーザー検知が確定した時刻を格納
  int correctedBPM; // 今回レーザーが通過したときの間隔（BPM）
  int stableCount; // 同じBPMが何回連続で出たか
  int prevCorrectedBPM; // 前回レーザーが通過したときの間隔(BPM)
  bool laserOn; // 現在レーザーがCdSセルに当たっているかを格納
  int baseline; // 暗状態（レーザー非通過時）の読み取り値の追従値を格納
  bool baselineReady; // baselineの初期値取得が完了したかを格納
  int onThreshold; // 現在有効なON判定しきい値を格納
  int offThreshold; // 現在有効なOFF判定しきい値を格納
  int pulsePeak; // 検知中のパルスにおける読み取り値の最大値を一時的に格納
  int peakHistory[CDS_PEAK_HISTORY]; // しきい値較正に用いる直近ピーク値を保持する配列
  uint8_t peakIndex; // peakHistory内の現在の書き込み位置を管理するインデックス
  uint8_t peakCount; // peakHistoryに記録済みのピーク数を格納
  bool calibrated; // しきい値の較正が一度でも行われたかを格納
  bool startupCalibrated; // 起動直後のキャリブレーション期間が終了したかを格納
  bool waitingForRelease; // 起動時に暗転待ち状態かを格納
  unsigned long calibrationStartMs; // 起動キャリブレーションを開始した時刻を格納
  int startupMin; // 起動キャリブレーション中に観測された最小値を格納
  int releaseThreshold; // waitingForRelease状態を解除する判定値を格納

  // レーザーのON/OFF判定，起動時キャリブレーション，ベースライン更新を行う．
  bool detectLight(int v, unsigned long nowMs) {
    if (!baselineReady) {
      baseline = v;
      startupMin = v;
      baselineReady = true;
      updateFallbackThresholds();
    }

    // 起動直後にレーザーが当たっていても、最初の1秒は検出せず暗い側の値を探す
    if (!startupCalibrated) {
      if (v < startupMin) startupMin = v;
      baseline = startupMin;
      updateFallbackThresholds();
      if (nowMs - calibrationStartMs < CDS_STARTUP_CALIBRATION_MS) return false;
      startupCalibrated = true;
      waitingForRelease = (v >= onThreshold);
      releaseThreshold = (baseline > 1023 - CDS_MIN_DELTA) ? baseline - CDS_MIN_DELTA : offThreshold;
      return false;
    }

    if (waitingForRelease) {
      if (v <= releaseThreshold) {
        baseline = v;
        updateFallbackThresholds();
        waitingForRelease = false;
      }
      return false;
    }

    if (!laserOn) {
      if (v < baseline) baseline = v;
      else baseline = (baseline * 31 + v) / 32;
      if (!calibrated) updateFallbackThresholds();
    }

    if (!laserOn && v >= onThreshold) {
      laserOn = true;
      pulsePeak = v;
      return true;
    }

    if (laserOn) {
      if (v > pulsePeak) pulsePeak = v;
      if (v <= offThreshold) {
        laserOn = false;
        registerPeak(pulsePeak);
      }
    }
    return false;
  }

  // 検知したピーク値を記録し，ON/OFF判定しきい値を再較正する．
  void registerPeak(int peak) {
    if (peak - baseline < CDS_MIN_DELTA) return;

    peakHistory[peakIndex] = peak;
    peakIndex = (peakIndex + 1) % CDS_PEAK_HISTORY;
    if (peakCount < CDS_PEAK_HISTORY) peakCount++;
    if (peakCount < 3) return;

    long sum = 0;
    for (uint8_t i = 0; i < peakCount; i++) sum += peakHistory[i];
    int peakAvg = sum / peakCount;
    int amplitude = peakAvg - baseline;
    if (amplitude < CDS_MIN_DELTA) return;

    int targetOn = constrain(baseline + (amplitude * CDS_ON_RATIO) / 100, 0, 1023);
    int targetOff = constrain(baseline + (amplitude * CDS_OFF_RATIO) / 100, 0, targetOn - 5);

    if (!calibrated) {
      onThreshold = targetOn;
      offThreshold = targetOff;
      calibrated = true;
    } else {
      onThreshold = limitStep(onThreshold, targetOn, CDS_THRESHOLD_STEP_LIMIT);
      offThreshold = limitStep(offThreshold, targetOff, CDS_THRESHOLD_STEP_LIMIT);
    }
    if (offThreshold >= onThreshold) offThreshold = max(0, onThreshold - 5);
  }

  // 較正前に使用する仮のしきい値を更新する．
  void updateFallbackThresholds() {
    onThreshold = constrain(baseline + CDS_MIN_DELTA, 0, 1023);
    offThreshold = constrain(baseline + CDS_MIN_DELTA / 2, 0, max(0, onThreshold - 5));
  }

  // しきい値が一度に変化する最大量を制限する．
  int limitStep(int current, int target, int limit) {
    if (target > current + limit) return current + limit;
    if (target < current - limit) return current - limit;
    return target;
  }

  // ピーク履歴を初期化する．
  void clearPeakHistory() {
    for (uint8_t i = 0; i < CDS_PEAK_HISTORY; i++) peakHistory[i] = 0;
  }

  // 実測の1周時間テーブルの中から最も近いものを探し、対応するBPM_LISTの値を返す
  int intervalToBpm(float intervalMs) {
    int best = BPM_LIST[0]; float bestD = fabs(intervalMs - ROTATION_PERIOD_LIST[0]);
    for (int i = 1; i < 5; i++) {
      float d = fabs(intervalMs - ROTATION_PERIOD_LIST[i]);
      if (d < bestD) { bestD = d; best = BPM_LIST[i]; }
    }
    return best;
  }
};

SyncManager syncManager;

// ==========================================
// state:
//   0 = 待機（BPM安定待ち）
//   1 = BPM安定、ドラムと同期中（Tパケット送信・diff収束待ち）
//   2 = 演奏中
//   4 = システム停止中
// ==========================================
int state = 4;


// 楽譜遷移（ジャンプ）定義
// jumpTable[曲ID][ジャンプ番号][0=トリガーtick, 1=ジャンプ先tick, 2=ジャンプ先scoreIdx]
const int MAX_SONGS = 5;
const int MAX_JUMPS = 20;
const int32_t jumpTable[MAX_SONGS][MAX_JUMPS][3] = {
  { // 曲ID: 0 かえるのうた
    {1536, 0, 0},
    {-1, -1, -1}
  },
  {
    {2880, 1728, 48}, {2784, 2880, 281}, {4512, 1728, 48},
    {-1, -1, -1}
  },
  {
    {5760, 4608, 651}, {5736, 5832, 850}, {6432, 5856, 856}, {6216, 6456, 993}, {6816, 4608, 651}, {5760, 6816, 1046}, {8832, 4608, 651},
    {-1, -1, -1}
  },
  {
    {9240, 8856, 1399}, {10008, 9240, 1499}, {10392, 10008, 1669}, {10776, 10392, 1783}, {11544, 10776, 1895}, {11544, 10008, 1669}, {10392, 10008, 1669}, {10392, 8856, 1399}, {9240, 8856, 1399}, {10008, 11544, 2111}, {11928, 11544, 2111}, {11904, 11952, 2225}, {13632, 8856, 1399},
    {-1, -1, -1}
  },
  {
    {16488, 14952, 2973},
    {-1, -1, -1}
  }
};

uint8_t currentJumpIdx = 0; // 現在のジャンプインデックス

// 曲ごとの初期設定データ
struct SongStartData {
  uint32_t startTick;
  uint16_t startNoteIdx;
};

// Processingの設定に合わせた初期値の定義
// インデックスがそのまま currentSongId に対応
const SongStartData songStartParams[MAX_SONGS] = {
  {0, 0},       // 曲ID: 0 (かえるのうた)
  {1536, 29},   // 曲ID: 1 (酒場でブギウギ)
  {4608, 651},  // 曲ID: 2 (王宮のメヌエット)
  {8856, 1399}, // 曲ID: 3 (トルコ行進曲)
  {13632, 2638} // 曲ID: 4 (序曲)
};

// ==========================================
// 関数宣言
// ==========================================
void updateTick();
void setBpm(uint8_t bpm);
void applyDiffCorrection(int32_t diff);
void sendTPacket();
void sendPacket();
void sendMPacket(uint8_t bpm);

// ==========================================
// セットアップ
// ==========================================
void setup() {
  Serial.begin(115200);
  WiFi.config(localIP);
  WiFi.begin(SSID, PASS);
  Serial.print("Connecting WiFi...");
  while (WiFi.status() != WL_CONNECTED) { delay(500); Serial.print("."); }
  Serial.println("\nWiFi Connected!");
  Udp.begin(LOCAL_PORT);
  Serial.println("state=0 BPM安定待ち");
}

// ==========================================
// メインループ
// ==========================================
void loop() {
  unsigned long now      = millis();
  int           cdsValue = analogRead(CDS_PIN);

  // ---- UDP受信 ----
  int pktSize = Udp.parsePacket();
  while (pktSize > 0) {
  if (pktSize >= 2) {
    Udp.read(packetBuffer, sizeof(packetBuffer));
    char    header  = (char)packetBuffer[0];
    uint8_t payload = packetBuffer[1];

    switch (header) {
      // 'R': ドラムからのdiff返信
      case 'R':
        // 停止中（4）や安定待ち（0）に届いた過去のR通信は遅延ゴミなので無視
        if (state == 0 || state == 4) break;

        if (pktSize >= 5) {
          int32_t rawDiff =
            ((int32_t)(int8_t)packetBuffer[1] << 24) |
            ((int32_t)packetBuffer[2] << 16) |
            ((int32_t)packetBuffer[3] <<  8) |
             (int32_t)packetBuffer[4];
          
          // 遅延（レイテンシ）補正の計算
          uint32_t rtt = micros() - lastTPacketMicroTime;
          // 【安全対策】往復時間が長すぎる場合は古いパケットとみなして破棄
          if (rtt > 1000 * 1000) {
            Serial.println(F("[R] パケット遅延過大のためスキップ"));
            break; 
          }

          uint32_t latencyMicros = rtt / 2; // 片道遅延
          int32_t latencyTicks = 0;

          if (baseTickIval > 0) {
            // 四捨五入してTick数に変換: (値 + 割る数の半分) / 割る数
            latencyTicks = (int32_t)((latencyMicros + (baseTickIval / 2)) / baseTickIval);
          }

          // ドラム機の勘違い（遅延分）を差し引いて真のズレを計算
          int32_t diff = rawDiff - latencyTicks; 

          Serial.print(F("[R] rawDiff=")); Serial.print(rawDiff);
          Serial.print(F(" latencyTicks=")); Serial.print(latencyTicks);
          Serial.print(F(" trueDiff=")); Serial.println(diff);

          // Serial.print(F("[R] diff=")); Serial.println(diff);

          if (drumNotStarted) {
            if ((int32_t)lastTPacketTick + diff == 0) {
              // Tパケット送信時点でドラムはまだ0。自身もリセットして足並みを揃える
              globalTick = 0;
              lastTick   = micros(); // Tick計測の基点もリセット
              Serial.println(F("[R] ドラム未起動: globalTick リセット"));
            } else {
              // ドラムが動き出した初回：即時同期
              drumNotStarted = false;
              globalTick = (uint32_t)((int32_t)globalTick + diff);
              lastDiff   = diff;
              Serial.print(F("[R] ドラム起動検知: 即時同期 globalTick="));
              Serial.println(globalTick);
              // 以降は通常補正も適用
              if (state >= 1) {
                // 「補正が完了している」または「緊急事態（ズレが24Tick以上）」の時だけ補正を更新
                if (ticksToCorrect == 0 || abs(diff) >= 24) {
                  applyDiffCorrection(diff);
                } else {
                  Serial.println(F("[R] 補正中のため新規R通信の適用をスキップ"));
                }
              }
            }
          } else {
            // 通常の同期補正
            lastDiff = diff;
            if (state >= 1) {
              // 「補正が完了している」または「緊急事態（ズレが24Tick以上）」の時だけ補正を更新
              if (ticksToCorrect == 0 || abs(diff) >= 24) {
                applyDiffCorrection(diff);
              } else {
                Serial.println(F("[R] 補正中のため新規R通信の適用をスキップ"));
              }
            }
          }
        }
        break;

      // 'O': ドラムからの演奏開始指示
      case 'O':
        if (pktSize >= 6) {
          if(waitingForStart || state == 2) return;
          
          uint8_t  songId  = packetBuffer[1];
          uint32_t sgt =
            ((uint32_t)packetBuffer[2] << 24) |
            ((uint32_t)packetBuffer[3] << 16) |
            ((uint32_t)packetBuffer[4] <<  8) |
             (uint32_t)packetBuffer[5];

          currentSongId   = songId;
          startGlobalTick = sgt;

          Serial.print(F("[O] songId=")); Serial.print(songId);
          Serial.print(F(" startGlobalTick=")); Serial.println(sgt);

          // 既に超過していたらスキップして即開始
          if ((int32_t)(globalTick - sgt) >= 0) {
            // 超過分だけ curTick を進める
            uint32_t overrun = globalTick - sgt;
            curTick   = songStartParams[currentSongId].startTick + overrun;
            scoreIdx  = songStartParams[currentSongId].startNoteIdx;
            // overrunより前のノートをスキップ
            while (scoreIdx < note_count && Note[scoreIdx].timing < curTick) scoreIdx++;
            state = 2;
            Serial.print(F("[O] 超過済み。overrun=")); Serial.print(overrun);
            Serial.println(F(" Tick → 即開始"));
          } else {
            waitingForStart = true;
            Serial.println(F("[O] 開始待機中..."));
          }
        }
        break;

      case 'V': {
        uint8_t vol = (payload >= 1 && payload <= 5) ? (payload * 10) : 30;
        if (vol != curVol) {
          curVol = vol;
          Serial.print('V'); Serial.print(','); Serial.print(vol); Serial.print('\n');
        }
        break;
      }

      // 'E': 演奏終了通知（指揮機から）
      case 'E':
        drumNotStarted  = true;
        waitingForStart = false;
        state           = 4;  // BPMは確定済みのままstate1に戻す 観覧車も停止するため1ではなく0
        curTick         = 0;
        scoreIdx        = 0;
        globalTick      = 0;
        curBpm          = 0;  // 追加: BPMリセット
        tickIval        = 0;  // 追加
        baseTickIval    = 0;  // 追加
        Serial.println(F("[E] 演奏終了。drumNotStarted リセット"));
        break;

      // -------------------------------------------------------
      // 'S': 演奏開始通知（指揮機から）．ドラムに対して演奏への参加を申請
      case 'S':
        drumNotStarted  = true;
        waitingForStart = false;
        state           = 0;  // BPMは確定済みのままstate1に戻す 観覧車も停止するため1ではなく0
        curTick         = 0;
        scoreIdx        = 0;
        globalTick      = 0;
        currentJumpIdx  = 0; // ジャンプインデックスも初期化
        Serial.println(F("[S] 演奏開始。drumNotStarted リセット"));

        syncManager.reset();

        sendJPacket();

        Serial.println(F("[S] ドラムに同期予告を送信。"));

        break;

        // マスターからの確定BPMを受信
      case 'N':
        if (pktSize >= 2) {
          uint8_t confirmedBpm = packetBuffer[1];
          setBpm(confirmedBpm);
          Serial.print(F("[N] 確定BPM受信: ")); Serial.println(confirmedBpm);
        }
        break;
      

      default: break;
    }
  } else {
    Udp.read(); // 空読みして破棄
  }
  pktSize = Udp.parsePacket(); // 次のパケットがあるか確認
  }

  if (state == 4) return;

  // ---- レーザー検知 ----
  bool laser = syncManager.update(cdsValue, now);

  // ---- ステートマシン ----
  switch (state) {

    // --- state0: BPM安定待ち ---
    case 0:
      if (laser) {
        int s = syncManager.getStableCount();
        uint8_t curLaserBpm = syncManager.getCorrectedBPM();
        Serial.print(F("[S0] BPM=")); Serial.print(curLaserBpm);
        Serial.print(F(" stable=")); Serial.println(s);

        sendMPacket(curLaserBpm); // M通信でドラムに提案

        if (s >= MIN_STABLE_COUNT) {
          state = 1;
          lastTick = micros();
          // setBpm(syncManager.getCorrectedBPM());

          Serial.println(F("[S0→1] BPM安定(N通信待ち)"));
        }
      }
      break;

    // --- state1: 同期中 ---
    case 1:
      updateTick();
      if (laser) {
        static unsigned long lastSentTime = 0;
        if (now - lastSentTime > 200) { // 200ms以内の連続送信（チャタリング）をブロック
        // setBpm(syncManager.getCorrectedBPM());
        // レーザー通過のたびにTパケットを送信してdiff返信を要求
          sendMPacket(syncManager.getCorrectedBPM()); // M通信で提案
          lastTPacketTick = globalTick;
          sendTPacket();
          lastSentTime = now;
        }
      }

      // waitingForStart かつ startGlobalTick に到達したら演奏開始
      if (waitingForStart && (int32_t)(globalTick - startGlobalTick) >= 0) {
        waitingForStart = false;
        curTick  = songStartParams[currentSongId].startTick;
        scoreIdx = songStartParams[currentSongId].startNoteIdx;
        state    = 2;
        Serial.println(F("[S1→2] 演奏開始！"));
      }
      break;

    // --- state2: 演奏中 ---
    case 2:
      updateTick();
      if (laser) {
        static unsigned long lastSentTime = 0;
        if (now - lastSentTime > 200) { //200ms以内の連続送信（チャタリング）をブロック
        // setBpm(syncManager.getCorrectedBPM());
        // レーザー通過のたびにTパケットを送信してdiff返信を要求
          sendMPacket(syncManager.getCorrectedBPM()); // M通信で提案
          lastTPacketTick = globalTick;
          sendTPacket();
          lastSentTime = now;
        }
      }
      break;
    
    // --- state4: システム停止中 ---
    // case 4: break; 
  }
}

// ==========================================
// diff補正によるtickIval算出
//   diff = drumGlobalTick - myGlobalTick
//   diff > 0: 自分が遅れている → 速く進む（tickIvalを小さく）
//   diff < 0: 自分が早すぎる  → 遅く進む（tickIvalを大きく）
// ==========================================
void applyDiffCorrection(int32_t diff) {
  if (baseTickIval == 0) return;

  float SYNC_READY_THRESHOLD = 60 * 1000 / baseTickIval;
  
  if (abs(diff) <= SYNC_READY_THRESHOLD) {
    // ズレが±2 Tick以内なら何もしない
    tickIval = baseTickIval;
    ticksToCorrect = 0;
  } else if (diff <= -24) {
    // 大幅に早い → ×2 で -diff Tick間遅らせる
    tickIval       = baseTickIval * 2;
    ticksToCorrect = -diff;
    Serial.println(F("[SYNC] 大幅早い: ×2"));
  } else if (diff >= 24) {
    // 大幅に遅い → ×2/3 で diff×3 Tick間速める
    tickIval       = baseTickIval * 2 / 3;
    ticksToCorrect = diff * 3;
    Serial.println(F("[SYNC] 大幅遅い: ×2/3"));
  } else {
    // 小さいズレ → 次の48Tickでdiffを吸収する比率補正
    tickIval       = (uint32_t)((48UL * baseTickIval) / (48 + diff));
    ticksToCorrect = 48 + diff;
    Serial.print(F("[SYNC] 微補正 diff=")); Serial.println(diff);
  }

  uint8_t sendProcessingBpm = (uint8_t)(60000000UL / (tickIval * 24));
  Serial.print('B'); Serial.print(','); Serial.print(sendProcessingBpm); Serial.print('\n');
}

// ==========================================
// BPM更新
// ==========================================
void setBpm(uint8_t bpm) {
  if (bpm == curBpm) return;
  curBpm       = bpm;
  uint32_t oldBaseTickIval = baseTickIval;
  baseTickIval = 60000000UL / (uint32_t)bpm / 24;

  if (oldBaseTickIval > 0 && ticksToCorrect > 0) {
    // 現在の補正スピードを新しいBPMの比率に合わせてスケールさせる
    tickIval = (tickIval * baseTickIval) / oldBaseTickIval;
    Serial.println(F("[BPM変化] 補正スピードを新BPMにスケール"));
  } else {
    tickIval = baseTickIval;
  }

  uint8_t sendProcessingBpm = (uint8_t)(60000000UL / (tickIval * 24));
  Serial.print('B'); Serial.print(','); Serial.print(sendProcessingBpm); Serial.print('\n');
}

// Tパケット送信（自身のglobalTickをドラムへ）
void sendTPacket() {
  Udp.beginPacket(drumIP, LOCAL_PORT);
  Udp.write('T');
  Udp.write((uint8_t)((globalTick >> 24) & 0xFF));
  Udp.write((uint8_t)((globalTick >> 16) & 0xFF));
  Udp.write((uint8_t)((globalTick >>  8) & 0xFF));
  Udp.write((uint8_t)( globalTick        & 0xFF));
  Udp.write(MY_INST_ID);
  Udp.endPacket();
  lastTPacketMicroTime = micros();
}

void sendJPacket() {
  for(int i=0;i<3;i++){
    Udp.beginPacket(drumIP, LOCAL_PORT);
    Udp.write('J');
    Udp.write(MY_INST_ID);
    Udp.endPacket();
    delay(3); // Wi-Fiのバッファ溢れを防ぐための必須の遅延
  }
}

// M通信送信（自身のレーザーBPMをドラムへ提案）
void sendMPacket(uint8_t bpm) {
    Udp.beginPacket(drumIP, LOCAL_PORT);
    Udp.write('M');
    Udp.write(bpm);
    Udp.write(MY_INST_ID);
    Udp.endPacket();
}

// ==========================================
// Tick管理
// ==========================================
void updateTick() {
  if (tickIval == 0) return;
  uint32_t now = micros();
  while (now - lastTick >= tickIval) {
    uint32_t currentIval = tickIval; // ループ突入時の値を保存

    globalTick++;

    // 補正カウントダウン
    if (ticksToCorrect > 0) {
      ticksToCorrect--;
      if (ticksToCorrect == 0) {
        tickIval = baseTickIval;
        Serial.println(F("[SYNC] 補正完了"));
      }
    }

    if (state == 2) {
      // ジャンプ（ループ）処理の判定
      if (jumpTable[currentSongId][currentJumpIdx][0] != -1 && curTick == jumpTable[currentSongId][currentJumpIdx][0]) {
        curTick  = jumpTable[currentSongId][currentJumpIdx][1];
        scoreIdx = jumpTable[currentSongId][currentJumpIdx][2];
        currentJumpIdx++;
        
        if (currentJumpIdx >= MAX_JUMPS || jumpTable[currentSongId][currentJumpIdx][0] == -1) {
          currentJumpIdx = 0;
        }
      }

      sendPacket();
      curTick++;
    }

    lastTick += currentIval; // ★書き換えられていても元の値を足す
  }
}

// ==========================================
// ノート送信（Processingへ）
// ==========================================
void sendPacket() {
  while (scoreIdx < note_count && Note[scoreIdx].timing <= curTick) {
    //uint8_t vs = (uint8_t)((uint16_t)Note[scoreIdx].startVelocity * curVol / 50);
    //uint8_t ve = (uint8_t)((uint16_t)Note[scoreIdx].endVelocity   * curVol / 50);
    Serial.print('N'); Serial.print(',');
    Serial.print(Note[scoreIdx].pitch);    Serial.print(',');
    Serial.print(Note[scoreIdx].duration); Serial.print(',');
    Serial.print(Note[scoreIdx].startVelocity); Serial.print(',');
    Serial.print(Note[scoreIdx].endVelocity); Serial.print(',');
    Serial.println(Note[scoreIdx].type);
    scoreIdx++;
  }
}
\end{lstlisting}

\subsection{strings\_0706.ino（スレーブ機・相違箇所のみ）}
\begin{lstlisting}[caption={strings\_0706.ino（piano\_0706.inoとの相違箇所のみ）}, label={lst:douki_strings_full}]
// strings_0706.ino：以下の2定数と楽譜データ（Note[]配列）を除き，
// piano_0706.ino と完全に同一である．

IPAddress localIP(192, 168, 11, 13);

// 楽器ID（0=ドラム, 1=ピアノ, 2=ストリングス, 3=フルート）
const uint8_t MY_INST_ID = 2; // ストリングス
\end{lstlisting}

\subsection{flute\_0706.ino（スレーブ機・相違箇所のみ）}
\begin{lstlisting}[caption={flute\_0706.ino（piano\_0706.inoとの相違箇所のみ）}, label={lst:douki_flute_full}]
// flute_0706.ino：以下の2定数と楽譜データ（Note[]配列）を除き，
// piano_0706.ino と完全に同一である．

IPAddress localIP(192, 168, 11, 14);

// 楽器ID（0=ドラム, 1=ピアノ, 2=ストリングス, 3=フルート）
const uint8_t MY_INST_ID = 3; // フルート
\end{lstlisting}
